% GENERATED FILE. Edit canonical lessons, then run npm run generate:books.
% canonical-source-hash: fnv1a64:0c33b754829ba319
% canonical-lessons: HI-C528-tarasna, HI-C528-naqqasi, HI-C528-silp, HI-C528-godna, HI-C528-pradarsani

\chapter{Craft and Carving}
\label{ch:hi-craft-and-carving}
\hlchaptermodality{528}

\begin{chapteropening}
\textbf{By the end of this chapter:} \emph{I can talk about carving, crafts and exhibitions.}

Complete the last lesson of chapter 528: I can talk about carving, crafts and exhibitions.
\end{chapteropening}

\section[\texorpdfstring{\hi{तराशना} (tarāśnā)}{तराशना (tarāśnā)}]{\hi{तराशना} (tarāśnā) — to carve; to cut into shape}

\label{lesson:HI-C528-tarasna}



\begin{quote}
Before the new one: say the Hindi for a line (drawn), then the Hindi for to colour.
\end{quote}

\subsection*{You'll want to know: \hi{तराशना}}
\textbf{\hi{तराशना}} — \emph{tarāśnā} — \textquotedblleft{}to carve; to cut into shape\textquotedblright{}.

\begin{quote}
\hi{कारीगर} \hi{ने} \hi{पत्थर} \hi{तराशकर} \hi{एक} \hi{हाथी} \hi{बनाया।}
\end{quote}

\begin{quote}
\emph{Kārīgar ne patthar tarāśkar ek hāthī banāyā.}
\end{quote}

\begin{quote}
The craftsman carved the stone into an elephant.
\end{quote}

Past, with \hi{ने} (he or she …): \hi{उसने} \hi{पत्थर} \hi{तराशा}, \emph{usne patthar tarāśā}.

\begin{cousinweb}
\textbf{\hi{तराशना}} (\emph{tarāśnā}) is usually traced to Persian \emph{tarāsh}, \textquotedblleft{}cutting\textquotedblright{}, the stem of \emph{tarāshīdan}, \textquotedblleft{}to shave, to carve\textquotedblright{}, with the Hindi verb ending -nā.
\end{cousinweb}

\begin{grammarlens}[title={What you've built}]
One more word for describing what you lived and felt.
\end{grammarlens}

\subsection*{Your turn}
\begin{itemize}
\raggedright
  \item \emph{Say it:} \emph{tarāśnā}
  \item \emph{Say it:} \emph{tarāśnā}, once more
  \item \emph{Say it:} \emph{Kārīgar ne patthar tarāśkar ek hāthī banāyā.} — The craftsman carved the stone into an elephant.
  \item \emph{Recall:} say the Hindi for a line (drawn), then the Hindi for to colour, then say \emph{tarāśnā} again
\end{itemize}

\begin{tcolorbox}[breakable,colback=teal!4,colframe=teal!35!black,title={Before you move on}]
What does \hi{तराशना} mean? (\textquotedblleft{}To carve; to cut into shape\textquotedblright{}.) Say it once more.
\end{tcolorbox}
\section[\texorpdfstring{\hi{नक़्क़ाशी} (naqqāśī)}{नक़्क़ाशी (naqqāśī)}]{\hi{नक़्क़ाशी} (naqqāśī) — carving; engraved decoration}

\label{lesson:HI-C528-naqqasi}



\begin{quote}
Before the new one: say the Hindi for to colour, then the Hindi for to carve.
\end{quote}

\subsection*{You'll want to know: \hi{नक़्क़ाशी}}
\textbf{\hi{नक़्क़ाशी}} — \emph{naqqāśī} — \textquotedblleft{}carving; engraved decoration\textquotedblright{}.

\begin{quote}
\hi{महल} \hi{के} \hi{दरवाज़ों} \hi{पर} \hi{सुंदर} \hi{नक़्क़ाशी} \hi{है।}
\end{quote}

\begin{quote}
\emph{Mahal ke darvāzoṁ par sundar naqqāśī hai.}
\end{quote}

\begin{quote}
There is beautiful carving on the palace doors.
\end{quote}

Feminine.

\begin{cousinweb}
\textbf{\hi{नक़्क़ाशी}} (\emph{naqqāśī}) is usually traced to Arabic \emph{naqqāsh}, \textquotedblleft{}a painter, an engraver\textquotedblright{}, through Persian. Hindi adds -ī for the work itself.
\end{cousinweb}

\begin{grammarlens}[title={What you've built}]
One more word for describing what you lived and felt.
\end{grammarlens}

\subsection*{Your turn}
\begin{itemize}
\raggedright
  \item \emph{Say it:} \emph{naqqāśī}
  \item \emph{Say it:} \emph{naqqāśī}, once more
  \item \emph{Say it:} \emph{Mahal ke darvāzoṁ par sundar naqqāśī hai.} — There is beautiful carving on the palace doors.
  \item \emph{Recall:} say the Hindi for to colour, then the Hindi for to carve, then say \emph{naqqāśī} again
\end{itemize}

\begin{tcolorbox}[breakable,colback=teal!4,colframe=teal!35!black,title={Before you move on}]
What does \hi{नक़्क़ाशी} mean? (\textquotedblleft{}Carving; engraved decoration\textquotedblright{}.) Say it once more.
\end{tcolorbox}
\section[\texorpdfstring{\hi{शिल्प} (śilp)}{शिल्प (śilp)}]{\hi{शिल्प} (śilp) — a craft; craftsmanship}

\label{lesson:HI-C528-silp}



\begin{quote}
Before the new one: say the Hindi for to carve, then the Hindi for carving.
\end{quote}

\subsection*{You'll want to know: \hi{शिल्प}}
\textbf{\hi{शिल्प}} — \emph{śilp} — \textquotedblleft{}a craft; craftsmanship\textquotedblright{}.

\begin{quote}
\hi{यह} \hi{गाँव} \hi{लकड़ी} \hi{के} \hi{शिल्प} \hi{के} \hi{लिए} \hi{मशहूर} \hi{है।}
\end{quote}

\begin{quote}
\emph{Yah gā̃v lakṛī ke śilp ke liye maśhūr hai.}
\end{quote}

\begin{quote}
This village is famous for its woodcraft.
\end{quote}

Masculine.

\begin{cousinweb}
\textbf{\hi{शिल्प}} (\emph{śilp}) is a tatsama, taken whole from Sanskrit \emph{śilpa}, \textquotedblleft{}artistic work, decoration\textquotedblright{} (Monier-Williams).
\end{cousinweb}

\begin{grammarlens}[title={What you've built}]
One more word for describing what you lived and felt.
\end{grammarlens}

\subsection*{Your turn}
\begin{itemize}
\raggedright
  \item \emph{Say it:} \emph{śilp}
  \item \emph{Say it:} \emph{śilp}, once more
  \item \emph{Say it:} \emph{Yah gā̃v lakṛī ke śilp ke liye maśhūr hai.} — This village is famous for its woodcraft.
  \item \emph{Recall:} say the Hindi for to carve, then the Hindi for carving, then say \emph{śilp} again
\end{itemize}

\begin{tcolorbox}[breakable,colback=teal!4,colframe=teal!35!black,title={Before you move on}]
What does \hi{शिल्प} mean? (\textquotedblleft{}A craft; craftsmanship\textquotedblright{}.) Say it once more.
\end{tcolorbox}
\section[\texorpdfstring{\hi{गोदना} (godnā)}{गोदना (godnā)}]{\hi{गोदना} (godnā) — to prick, to tattoo}

\label{lesson:HI-C528-godna}



\begin{quote}
Before the new one: say the Hindi for carving, then the Hindi for a craft.
\end{quote}

\subsection*{You'll want to know: \hi{गोदना}}
\textbf{\hi{गोदना}} — \emph{godnā} — \textquotedblleft{}to prick, to tattoo\textquotedblright{}.

\begin{quote}
\hi{मेले} \hi{में} \hi{एक} \hi{औरत} \hi{लोगों} \hi{के} \hi{हाथ} \hi{पर} \hi{फूल} \hi{गोद} \hi{रही} \hi{थी।}
\end{quote}

\begin{quote}
\emph{Mele meṁ ek aurat logoṁ ke hāth par phūl god rahī thī.}
\end{quote}

\begin{quote}
At the fair a woman was tattooing flowers on people's hands.
\end{quote}

Past, with \hi{ने} (he or she …): \hi{उसने} \hi{हाथ} \hi{पर} \hi{नाम} \hi{गोदा}, \emph{usne hāth par nām godā}.

\begin{cousinweb}
\textbf{\hi{गोदना}} (\emph{godnā}) is a tadbhava, inherited from \emph{*khōdd}, \textquotedblleft{}dig\textquotedblright{}, a form CDIAL reconstructs from its descendants (CDIAL 3934). CDIAL gives the Hindi verb as \textquotedblleft{}to pierce, hoe\textquotedblright{}.
\end{cousinweb}

\begin{grammarlens}[title={What you've built}]
One more word for describing what you lived and felt.
\end{grammarlens}

\subsection*{Your turn}
\begin{itemize}
\raggedright
  \item \emph{Say it:} \emph{godnā}
  \item \emph{Say it:} \emph{godnā}, once more
  \item \emph{Say it:} \emph{Mele meṁ ek aurat logoṁ ke hāth par phūl god rahī thī.} — At the fair a woman was tattooing flowers on people's hands.
  \item \emph{Recall:} say the Hindi for carving, then the Hindi for a craft, then say \emph{godnā} again
\end{itemize}

\begin{tcolorbox}[breakable,colback=teal!4,colframe=teal!35!black,title={Before you move on}]
What does \hi{गोदना} mean? (\textquotedblleft{}To prick, to tattoo\textquotedblright{}.) Say it once more.
\end{tcolorbox}
\section[\texorpdfstring{\hi{प्रदर्शनी} (pradarśanī)}{प्रदर्शनी (pradarśanī)}]{\hi{प्रदर्शनी} (pradarśanī) — an exhibition}

\label{lesson:HI-C528-pradarsani}



\begin{quote}
Before the new one: say the Hindi for a craft, then the Hindi for to prick, to tattoo.
\end{quote}

\subsection*{You'll want to know: \hi{प्रदर्शनी}}
\textbf{\hi{प्रदर्शनी}} — \emph{pradarśanī} — \textquotedblleft{}an exhibition\textquotedblright{}.

\begin{quote}
\hi{संग्रहालय} \hi{में} \hi{चित्रों} \hi{की} \hi{एक} \hi{प्रदर्शनी} \hi{लगी} \hi{है।}
\end{quote}

\begin{quote}
\emph{Saṅgrahālay meṁ citroṁ kī ek pradarśanī lagī hai.}
\end{quote}

\begin{quote}
There is an exhibition of paintings at the museum.
\end{quote}

Feminine.

\begin{cousinweb}
\textbf{\hi{प्रदर्शनी}} (\emph{pradarśanī}) is a modern formation built on Sanskrit \emph{pradarśana}, \textquotedblleft{}look, appearance; showing\textquotedblright{} (Monier-Williams), with the feminine -ī: a place where things are shown.
\end{cousinweb}

\begin{grammarlens}[title={What you've built}]
One more word for describing what you lived and felt.
\end{grammarlens}

\subsection*{Your turn}
\begin{itemize}
\raggedright
  \item \emph{Say it:} \emph{pradarśanī}
  \item \emph{Say it:} \emph{pradarśanī}, once more
  \item \emph{Say it:} \emph{Saṅgrahālay meṁ citroṁ kī ek pradarśanī lagī hai.} — There is an exhibition of paintings at the museum.
  \item \emph{Recall:} say the Hindi for a craft, then the Hindi for to prick, to tattoo, then say \emph{pradarśanī} again
\end{itemize}

\begin{tcolorbox}[breakable,colback=teal!4,colframe=teal!35!black,title={Before you move on}]
What does \hi{प्रदर्शनी} mean? (\textquotedblleft{}An exhibition\textquotedblright{}.) Say it once more.
\end{tcolorbox}
